% Compressed Appendix A; the expanded source is retained with the release.
% Each consolidated display retains one semantic label; external references are preserved.
Throughout, task coordinates use the common normalization of Section~\ref{sec:geometry}, and environments follow the declared distribution $\Pi$.

\subsection{Projection, Normalization, and Task--Nuisance Tangency}
\label{app:theory-projection}

Use the centered, square-integrable scores and blocks $J_{AB,e}$ of Section~\ref{sec:task-information}. Suppose the $d_N$ coordinates of the finite nuisance-score vector span $\mathcal T_{N,e}$ and $J_{NN,e}\succ0$. With $\Gamma_e=J_{TN,e}J_{NN,e}^{-1}$,
\begin{equation}
\begin{aligned}
\mathbb E_{P_e}[(s_{T,e}-\Gamma_es_{N,e})s_{N,e}^{\top}]&=J_{TN,e}-\Gamma_eJ_{NN,e}=0,\\
F_e=\mathbb E_{P_e}[s_{T,e}^{\eff}(s_{T,e}^{\eff})^{\top}]&=J_{TT,e}-J_{TN,e}J_{NN,e}^{-1}J_{NT,e}.
\end{aligned}
\label{eq:app-efficient-schur}
\end{equation}
Because $\Gamma_es_{N,e}$ lies componentwise in the nuisance span, orthogonality identifies the residual as the efficient score in Equation~\eqref{eq:efficient-score-information} \citep{bickel1993efficient,tsiatis2006semiparametric}. Equivalently,
\begin{equation}
u^{\top}F_eu=\min_{v\in\mathbb R^{d_N}}\mathbb E_{P_e}[(u^{\top}s_{T,e}-v^{\top}s_{N,e})^2],
\qquad v^\star=J_{NN,e}^{-1}J_{NT,e}u.
\label{eq:app-variational-schur}
\end{equation}
Hence $F_e\succeq0$, with $u^{\top}F_eu=0$ exactly when $u^{\top}s_{T,e}$ belongs to the nuisance span. The projection definition extends directly to the closed span of an infinite-dimensional nuisance tangent space.

For separate invertible transformations $\widetilde s_T=Us_T$ and $\widetilde s_N=Vs_N$, the nuisance span is unchanged and
\begin{equation}
\widetilde s_T^{\eff}=Us_T^{\eff},\qquad \widetilde F=UFU^{\top};
\qquad \frac{F}{J_{TT}}=1-\rho^2\quad\text{in the scalar case}.
\label{eq:app-coordinate-congruence}
\end{equation}
Here $\rho=J_{TN}/\sqrt{J_{TT}J_{NN}}$ requires positive raw task and nuisance information. Exact information is invariant to nuisance-basis changes; task congruence preserves rank but generally changes eigenvalues. The scalar surviving fraction is invariant to nonzero task rescaling and invertible nuisance-basis changes. Nuisance ridge preserves the latter invariance only if $\lambda_{\rid}I$ becomes $\lambda_{\rid}VV^\top$ under $s_N\mapsto Vs_N$ (Appendix~\ref{app:estimation-certification}).

For the tangency construction in Section~\ref{sec:tangency-separation}, let $b_e,r_e\in L_0^2(P_e)$ be orthogonal unit vectors, $\mathcal T_{N,e}=\operatorname{span}\{b_e\}$, and $a_e=b_e+\delta_er_e$ with $\delta_e\ge0$. The projection coefficient is $\langle a_e,b_e\rangle_e/\|b_e\|_e^2=1$, so
\begin{equation}
a_e^{\eff}=\delta_er_e,\qquad F_e=\|\delta_er_e\|_e^2=\delta_e^2,
\qquad J_{TT,e}=1+\delta_e^2,\qquad \frac{F_e}{J_{TT,e}}=\frac{\delta_e^2}{1+\delta_e^2}.
\label{eq:app-tangency-identities}
\end{equation}
This establishes Equation~\eqref{eq:tangency-identities}; coexistence with low prediction risk additionally requires a successful nuisance-based predictor.

\subsection{Normal Coordinates and the Inverse-Information Tail}
\label{app:theory-tails}

Under Theorem~\ref{thm:inverse-information-geometry}, the zero set $\mathcal Z=\{\lambda=0\}$ is a nonempty compact $C^2$ embedded manifold without boundary, of codimension $c\in\{1,\ldots,d\}$. Write $U_r(\mathcal Z)=\{z:\operatorname{dist}(z,\mathcal Z)<r\}$ for its tube, $N_s\mathcal Z$ for the normal space at $s$, $S(N_s\mathcal Z)$ for its unit sphere, and $\operatorname{vol}_{\mathcal Z}$ for induced volume. The nonnegative $\lambda$ is continuous on $\overline U_{r_\star}(\mathcal Z)$ and bounded away from zero outside it. In coordinates $z=s+r\omega$, assume uniformly over $\omega\in S(N_s\mathcal Z)$ and $s\in\mathcal Z$
\begin{equation}
\lambda(s+r\omega)=a(s,\omega)r^{2m}(1+o(1)),\qquad
p(s+r\omega)=b(s,\omega)r^\nu(1+o(1)),
\label{eq:null-set-expansions}
\end{equation}
where $p$ is the Lebesgue density, $m>0$, $\nu>-c$, and continuous $a,b$ have positive uniform lower and upper bounds. Put
\begin{equation}
\beta=\frac{c+\nu}{2m},\qquad
C_{\mathcal Z}=\frac1{c+\nu}\int_{\mathcal Z}\int_{S(N_s\mathcal Z)}
b(s,\omega)a(s,\omega)^{-\beta}\,d\sigma_s(\omega)\,d\operatorname{vol}_{\mathcal Z}(s).
\label{eq:app-tail-constant}
\end{equation}
The sphere measure $\sigma_s$ is unnormalized; for $c=1$, each of the two normal directions has mass one.

\begin{proof}[Proof of Theorem~\ref{thm:inverse-information-geometry}]
For any $0<r_0<r_\star$, continuity and absence of zeros on the compact annulus $r_0\le\operatorname{dist}(z,\mathcal Z)\le r_\star$ give a positive minimum of $\lambda$. Together with the bound outside the outer tube, this localizes $\{\lambda<t\}$ to $U_{r_0}(\mathcal Z)$ for all sufficiently small $t$.
The tubular volume element is $dz=j(s,r\omega)r^{c-1}\,dr\,d\sigma_s\,d\operatorname{vol}_{\mathcal Z}$, where the coordinate Jacobian satisfies $j=1+O(r)$ uniformly. Thus, for any $0<\zeta<1$, uniformity permits $r_0$ small enough that
\begin{equation}
\begin{aligned}
(1-\zeta)ar^{2m}&\le\lambda(s+r\omega)\le(1+\zeta)ar^{2m},\\
(1-\zeta)br^\nu&\le p(s+r\omega)j(s,r\omega)\le(1+\zeta)br^\nu
\end{aligned}
\qquad(0<r<r_0).
\label{eq:app-uniform-tube-bounds}
\end{equation}
For small $t$, all radial thresholds lie inside this tube. The event contains $r<(t/((1+\zeta)a))^{1/(2m)}$ and is contained in $r<(t/((1-\zeta)a))^{1/(2m)}$. Integrating uses $\int_0^{(t/a)^{1/(2m)}}br^{c+\nu-1}\,dr=ba^{-\beta}t^\beta/(c+\nu)$ and gives
\begin{equation}
(1-\zeta)(1+\zeta)^{-\beta}C_{\mathcal Z}
\le\frac{\Pr(\lambda<t)}{t^\beta}
\le(1+\zeta)(1-\zeta)^{-\beta}C_{\mathcal Z}.
\label{eq:app-small-information-sandwich}
\end{equation}
Let $t\downarrow0$ and then $\zeta\downarrow0$. Compactness and the coefficient bounds ensure $0<C_{\mathcal Z}<\infty$, so $\Pr(\lambda<t)\sim C_{\mathcal Z}t^\beta$. Since $c\ge1$, $\mathcal Z$ has zero Lebesgue and environment measure; substituting $t=x^{-1}$ proves the inverse tail in the next display.
\end{proof}

The tail-integral identity then yields, for every $p_0>0$,
\begin{equation}
\begin{aligned}
\Pr(\lambda^{-1}>x)&\sim C_{\mathcal Z}x^{-\beta},\\
\mathbb E[\lambda^{-p_0}]&=p_0\int_0^\infty x^{p_0-1}\Pr(\lambda^{-1}>x)\,dx<\infty
\quad\Longleftrightarrow\quad p_0<\beta.
\end{aligned}
\label{eq:app-inverse-moments}
\end{equation}
The bounded-$x$ part is finite; the tail has the convergence criterion of $\int^\infty x^{p_0-\beta-1}\,dx$. In particular, the inverse mean is finite exactly for $\beta>1$, with logarithmic divergence at $\beta=1$.

\paragraph{Scalar branches.}
For $\lambda=\delta^2$, $\Pr(\delta<t)\sim C_\delta t^\gamma$ gives exponent $\gamma/2$. More generally, let $\lambda_K=\sum_{k=1}^K w_k\delta_k^2$, $w_k>0$, where $\delta\in[0,\infty)^K$ has joint density $f$ continuous and positive at the origin from within the orthant. Changing variables gives the scalar portfolio law \mbox{\citep{mansouri_arai_2026_fisher_glasses}}:
\begin{equation}
\begin{aligned}
\Pr(\lambda_K<t)
&=\frac{t^{K/2}}{\sqrt{\prod_k w_k}}\int_{B_K^+}
 f\bigl(\sqrt{t/w_1}u_1,\ldots,\sqrt{t/w_K}u_K\bigr)\,du\\
&\sim\frac{f(0)v_K}{2^K\sqrt{\prod_k w_k}}\,t^{K/2}.
\end{aligned}
\label{eq:app-independent-branches}
\end{equation}
Here $B_K^+$ is the positive-orthant unit ball and $v_K$ the full ball's volume. Independence gives $f(0)=\prod_k f_k(0+)$ for branch marginal densities $f_k$; the inverse mean is finite for $K\ge3$. This calculation concerns additive scalar information or a fixed task direction. A minimum-eigenvalue conclusion instead requires applying Theorem~\ref{thm:inverse-information-geometry} to that eigenvalue.

\subsection{Local Classification: Testing Reduction and Tail Aggregation}
\label{app:theory-classification}

Let $\mathcal E_P(g)=R_P(g)-R_P^\star$. Fix a local family $\mathcal P_e^{\mathrm{loc}}$ independently of sample size, and define
\begin{equation}
\mathfrak R_e(n)=\inf_{\widehat g_n}\sup_{P\in\mathcal P_e^{\mathrm{loc}}}
\mathbb E_{S_n\sim P^{\otimes n}}[\mathcal E_P(\widehat g_n(S_n))],
\label{eq:app-local-minimax}
\end{equation}
where learning rules are measurable and $n$ counts independent observations from environment $e$; independent learner randomization is included in the expectation. Assume $\lambda_e>0$ and a unit weakest efficient-information direction $u_e$. For every $0<t\le t_{0,e}$, the fixed family contains $P_{e,\pm t}$ with task coordinates $\pm tu_e$ and potentially varying nuisance coordinates. For fixed $K_e,c_e,t_{0,e}>0$ and $q>0$, assume both bounds in Equation~\eqref{eq:local-classification-assumptions} hold throughout that radius.

\begin{proof}[Proof of Theorem~\ref{thm:local-classification-difficulty}]
Fix an admissible $t$ and abbreviate $P_\pm=P_{e,\pm t}$. For any learning rule, the decision-separation condition $\mathcal E_{P_+}(g)+\mathcal E_{P_-}(g)\ge c_et^q$ implies that choosing the alternative with smaller excess risk makes any testing error cost at least $c_et^q/2$ under the true alternative. Explicitly,
\begin{equation}
\begin{aligned}
\psi_t(S_n)=+&\ \Longleftrightarrow\ 
\mathcal E_{P_+}(\widehat g_n(S_n))\le\mathcal E_{P_-}(\widehat g_n(S_n)),\\
\max_{\sigma\in\{+,-\}}\mathbb E_{P_\sigma^{\otimes n}}[\mathcal E_{P_\sigma}(\widehat g_n)]
&\ge\frac{c_et^q}{4}\{P_+^{\otimes n}(\psi_t=-)+P_-^{\otimes n}(\psi_t=+)\}.
\end{aligned}
\label{eq:app-risk-to-testing}
\end{equation}
The standard exponential two-point testing inequality \citep[Lemma~2.6]{tsybakov2009introduction}, together with product additivity of KL and the assumed KL bound, gives
\begin{equation}
\begin{aligned}
P_+^{\otimes n}(\psi_t=-)+P_-^{\otimes n}(\psi_t=+)
&\ge\tfrac12\exp[-n\operatorname{KL}(P_+\Vert P_-)]
\ge\tfrac12 e^{-nK_et^2\lambda_e},\\
\mathfrak R_e(n)&\ge\frac{c_e}{8}t^q e^{-nK_et^2\lambda_e}.
\end{aligned}
\label{eq:app-testing-reduction}
\end{equation}
The second line takes the supremum over the fixed family containing both alternatives and then the infimum over learning rules. Choose
\begin{equation}
\begin{aligned}
t_n&=\min\{t_{0,e},(nK_e\lambda_e)^{-1/2}\},\qquad nK_e\lambda_et_n^2\le1,\\
\mathfrak R_e(n)&\ge c_0c_e\min\{t_{0,e}^{q},(nK_e\lambda_e)^{-q/2}\},
\qquad c_0=\frac{e^{-1}}8.
\end{aligned}
\label{eq:app-local-risk-bound}
\end{equation}
Independent learner randomization has the same law under both alternatives and leaves their KL divergence unchanged, so the argument also covers randomized rules.
\end{proof}

\paragraph{Quadratic decision separation.}
Suppose the binary alternatives share a feature marginal $Q_e$, with regression functions $\eta_{\pm,t}(x)=P_{e,\pm t}(Y=1\mid X=x)$ and Bayes rules $g_{\pm,t}^\star(x)=\mathbf1\{\eta_{\pm,t}(x)\ge1/2\}$. Let $D_{e,t}=\{g_{+,t}^\star\ne g_{-,t}^\star\}$ be their disagreement set. Applying the binary excess-risk identity \citep[Eq.~(1.1)]{audibert2007fast} gives
\begin{equation}
\begin{aligned}
&\inf_g\{\mathcal E_{P_{e,+t}}(g)+\mathcal E_{P_{e,-t}}(g)\}\\
&\qquad=\int_{D_{e,t}}\min\{|2\eta_{+,t}(x)-1|,|2\eta_{-,t}(x)-1|\}\,dQ_e(x).
\end{aligned}
\label{eq:app-boundary-separation}
\end{equation}
Pointwise, a classifier must disagree with at least one Bayes rule on $D_{e,t}$; choosing the smaller margin attains the integral for unrestricted measurable binary classifiers, while it remains a lower bound for a restricted class. If, for every $0<t\le t_{0,e}$, a measurable $A_{e,t}\subseteq D_{e,t}$ has $Q_e(A_{e,t})\ge c_{\mathrm{mass},e}t$ and both margins at least $c_{\mathrm{margin},e}t$, with positive constants independent of $t$, the integral is at least $c_{\mathrm{mass},e}c_{\mathrm{margin},e}t^2$. Thus $q=2$ and $c_e=c_{\mathrm{mass},e}c_{\mathrm{margin},e}$ satisfy decision separation; the KL condition is still required separately.

\paragraph{Environmental aggregation.}
If these assumptions hold for $\Pi$-almost every environment and all quantities are measurable, monotonicity and positive homogeneity of CVaR give
\begin{equation}
\begin{aligned}
\mathcal B_e(n)&=c_e\min\{t_{0,e}^{q},(nK_e\lambda_e)^{-q/2}\},\\
\operatorname{CVaR}_{\alpha,e\sim\Pi}(\mathfrak R_e(n))
&\ge c_0\operatorname{CVaR}_{\alpha,e\sim\Pi}(\mathcal B_e(n)).
\end{aligned}
\label{eq:app-tail-risk-aggregation}
\end{equation}
Separation $c_e$, distinguishability $K_e$, and radius $t_{0,e}$ remain inside the aggregation. For $q=2$ and $nK_e\lambda_e\ge t_{0,e}^{-2}$, $\mathcal B_e(n)=c_e/(nK_e\lambda_e)$; outside that interior regime the radius truncates the bound.

\subsection{CVaR, Information Coverage, and Stabilization}
\label{app:theory-certificate}

For $0<\alpha<1$ and nonnegative $Z$, the CVaR minimization formula holds for integrable $Z$ \citep[Theorem~10]{rockafellar2002conditional}; the threshold-weight argument below gives its dual:
\begin{equation}
\begin{aligned}
\operatorname{CVaR}_\alpha(Z)
&=\inf_{\tau\in\mathbb R}\left\{\tau+\frac{\mathbb E[(Z-\tau)_+]}{1-\alpha}\right\}
=\sup_{\xi\in\mathcal W_\alpha}\mathbb E[\xi Z],\\
\mathcal W_\alpha&=\{\xi:0\le\xi\le(1-\alpha)^{-1},\ \mathbb E\xi=1\}.
\end{aligned}
\label{eq:app-cvar-dual}
\end{equation}
Both values are $+\infty$ if $\mathbb EZ=+\infty$. For integrable $Z$, a maximizing weight equals $(1-\alpha)^{-1}$ above $v_\alpha=\inf\{v:\Pr(Z\le v)\ge\alpha\}$, zero below it, and fills the remaining mass at the threshold. Such fractional weighting is possible because $\Pr(Z>v_\alpha)\le1-\alpha\le\Pr(Z\ge v_\alpha)$. The dual gives monotonicity, positive homogeneity, and $v_\alpha\le\operatorname{CVaR}_\alpha(Z)$. For the finite audit distribution $\Pi_G$, let $\pi_g=\Pi_G(\{g\})\ge0$, $\sum_g\pi_g=1$. Then
\begin{equation}
\operatorname{CVaR}_{\alpha,\pi}(Z_g)
=\max_{\substack{\sum_g\omega_g=1\\0\le\omega_g\le\pi_g/(1-\alpha)}}\sum_g\omega_gZ_g.
\label{eq:app-finite-cvar}
\end{equation}
These normalized tail weights, including threshold atoms, also govern Appendix~\ref{app:tail-transverse-influence}.

\begin{proof}[Proof of Proposition~\ref{prop:information-coverage}]
Let $\varepsilon>0$ and $\lambda_e\in[0,\infty)$ be finite almost surely. Write $L_e^{(\varepsilon)}=(\lambda_e+\varepsilon)^{-1}$ and $I_\alpha^{\FG}(\varepsilon)$ for the certificate in Equation~\eqref{eq:fisher-glass-certificate}. If $I_\alpha^{\FG}(\varepsilon)\ge\rho>\varepsilon$, then
\begin{equation}
\begin{aligned}
v_\alpha(L^{(\varepsilon)})&\le\operatorname{CVaR}_\alpha(L^{(\varepsilon)})
=[I_\alpha^{\FG}(\varepsilon)]^{-1}\le\rho^{-1},\\
\Pr_{e\sim\Pi}(\lambda_e\ge\rho-\varepsilon)
&=\Pr_{e\sim\Pi}(L_e^{(\varepsilon)}\le\rho^{-1})\ge\alpha.
\end{aligned}
\label{eq:app-information-coverage}
\end{equation}
The last step is the quantile definition and includes equality at the information threshold.
\end{proof}
Since $0<L^{(\varepsilon)}\le\varepsilon^{-1}$, always $I_\alpha^{\FG}(\varepsilon)\ge\varepsilon$; only $\rho>\varepsilon$ certifies a strictly positive floor. With four equally weighted groups and $\alpha=0.95$, one group already contains the entire upper-tail mass, so
\begin{equation}
\operatorname{CVaR}_{0.95}(L^{(\varepsilon)})=\max_gL_g^{(\varepsilon)},
\qquad I_{0.95}^{\FG}(\varepsilon)=\min_g\lambda_g+\varepsilon.
\label{eq:app-four-group-certificate}
\end{equation}

\paragraph{Exact limit and positivity.}
Set $L^{(0)}=\lambda^{-1}$, with $1/0=+\infty$. As $\varepsilon\downarrow0$, $L^{(\varepsilon)}\uparrow L^{(0)}$, and monotone convergence applies to $\mathbb E[\xi L^{(\varepsilon)}]$ for every feasible $\xi$. The supremum over $\xi$ commutes with the increasing supremum along any sequence $\varepsilon_k\downarrow0$, giving
\begin{equation}
\begin{aligned}
\operatorname{CVaR}_\alpha(L^{(\varepsilon)})&\uparrow\operatorname{CVaR}_\alpha(\lambda^{-1}),\\
I_\alpha^{\FG}(\varepsilon)&\downarrow I_\alpha^{\FG}(0)
:=\bigl[\operatorname{CVaR}_\alpha(\lambda^{-1})\bigr]^{-1},\qquad 1/(+\infty)=0.
\end{aligned}
\label{eq:app-exact-limit}
\end{equation}
For nonnegative $Z$, the feasible weight $\xi=1$ and variational choice $\tau=0$ imply
\begin{equation}
\mathbb EZ\le\operatorname{CVaR}_\alpha(Z)\le\frac{\mathbb EZ}{1-\alpha},
\qquad I_\alpha^{\FG}(0)>0\ \Longleftrightarrow\ \mathbb E[\lambda^{-1}]<\infty.
\label{eq:app-cvar-mean-bounds}
\end{equation}
Under the geometric tail assumptions, this is equivalent to $\beta=(c+\nu)/(2m)>1$; the exact certificate vanishes for $\beta\le1$. Information coverage is a population statement under $\Pi$, separate from the confidence level of an estimated lower certificate.
